\paragraph{Linear coupling to the localized mode.}
For radiation of the same scalar field, a conditional projection identity
makes the distinction between scattering and local forcing precise.
In a fixed real $l=1$ harmonic, complexify the doubled reduced perturbation
$\psi=(rz_l,r\overline z_l)^T$ and write an additive source as $S$.
Equation~\eqref{nl:linear} and its conjugate give
\begin{equation}
 \psi_{tt}+2i\omega\sigma_3\psi_t+\mathcal K\psi=S,
 \qquad
 \mathcal K=-\partial_r^2+\frac{2}{r^2}
 +(D-\omega^2)I+C\sigma_1,
 \label{env:time}
\end{equation}
where $D,C$ are those of Eqs.~\eqref{eq:radialA}--\eqref{eq:radialB},
$\sigma_3=\operatorname{diag}(1,-1)$, and
$\sigma_1=\left(\begin{smallmatrix}0&1\\1&0\end{smallmatrix}\right)$.
The physical perturbation is recovered by the conjugate frequency pairing.
For a physical additive source in this real harmonic block,
$S_2=\overline{S_1}$; independent components are used only in the formal
complexification, with the conjugate pairing restored for physical forcing.
Use the radial inner product with measure $dr$ and the regular
self-adjoint domain of $\mathcal K$. With $X=(\psi,\psi_t)^T$, define
\begin{equation}
 \mathcal G=\begin{pmatrix}0&I\\-\mathcal K&-2i\omega\sigma_3\end{pmatrix},
 \qquad
 \mathcal J=\begin{pmatrix}2\omega\sigma_3&-iI\\iI&0\end{pmatrix}.
 \label{env:generator}
\end{equation}
Then $\mathcal G^\dagger\mathcal J+\mathcal J\mathcal G=0$ on that domain.
For a nonzero localized fundamental mode
$H(\rho)u=0$, where $H(\rho)=\mathcal K-\rho^2I-2\omega\rho\sigma_3$,
put $X_u=(u,i\rho u)^T$. If $\rho>|\omega|$, its conserved form is positive:
\begin{equation}
 N=\langle X_u,\mathcal JX_u\rangle
   =-\langle u,H'(\rho)u\rangle
   =2\langle u,(\rho I+\omega\sigma_3)u\rangle>0.
 \label{env:norm}
\end{equation}
Here $\rho$ is the fundamental frequency, not $2\rho$.
Positivity on this mode is neither a linear nor a nonlinear stability
theorem for the full system.
The corresponding left projection $c=\langle X_u,\mathcal JX\rangle/N$
obeys
\begin{equation}
 \dot c=i\rho c-\frac{i}{N}\langle u,S\rangle,
 \label{env:projection}
\end{equation}
provided the radial Green boundary form vanishes. On $[0,R]$ the
additional term is $i[u'^\dagger\psi-u^\dagger\psi']_0^R/N$.
At the origin this follows from regular values $u,\psi=O(r^2)$ and
regular derivatives $u',\psi'=O(r)$. At infinity the localized mode and
its derivative must decay against the controlled scattering field and
its derivative; this includes suitable finite wave packets.

Thus a homogeneous linear solution ($S=0$) with $c(0)=0$ cannot populate
this projection on the fixed background. Spatial separation of an
incoming packet alone does not imply exact orthogonality, because the
localized mode has tails. A prescribed local additive source can instead
have $\langle u,S\rangle\ne0$ and drive the mode. For $c(0)=0$ and a
finite square-integrable projected source, Eq.~\eqref{env:projection} gives
\begin{equation}
 N|c(T)|^2\leq\frac{T}{N}\int_0^T|\langle u,S(t)\rangle|^2\,dt.
 \label{env:pulse}
\end{equation}
This is a mode-projection bound, not a statement about the energy supplied
by an unspecified environment. Although Eq.~\eqref{env:norm} is positive
on this eigenspace, $\mathcal J$ is indefinite on the full phase space.
No globally passive port model, loading rate, nonlinear formation
mechanism, or indefinite maintenance follows. The local source and its
physical coupling must be specified separately; the nonlinear radiation
obstruction in Section~\ref{nl:section} addresses a different question.
